\scanpage{308}\mspage{32}
Aux fins de \unsure{référence|références} \unsure{ultérieure|ultérieures}, je vais
formuler maintenant les divers\deleted{\unsure{es}} aspects principaux
de l'\uline{axiome des} \inserted{$W$-}\uline{fibrations} $\mathrm{Loc}(7)$. La formulation
initiale à laquelle j'étais parvenu
(cf $\overline{\underline{\mathrm I}}$, $\overline{\underline{\mathrm{VI}}}$, $\overline{\underline{\mathrm{XII}}}$) \deleted{\unsure{du}} utilisait la notion
\unsure{que} j'appelle à présent ``$W$-fibration \uline{faible}'' \deleted{\unsure{,}}
(\inserted{l'}ex-$W$-fibration ; désormais détrônée !), et
sous sa forme la plus forte s'écrivait comme
l'inclusion
\[
\mathrm{Fib}f_W\subset\mathrm{Fib}_W\;.
\]
Cette notion de $W$-fibration faible, tout comme
celle de morphisme $W$-parfait, si \unsure{utile|utiles}
\unsure{soient-elles|soient-elle}, m'apparaît cependant un
peu technique, et je préfère donner
une formulation qui ne fasse appel
qu'à des notions plus familières et
d'apparence plus élémentaire. C'est pourquoi
je \deleted{\unsure{préfère me}} vais remplacer $\mathrm{Fib}f_W$ par
la sous-classe $\mathcal F_W\unsure{,}=\deleted{\unsure{\mathrm{Fib}}}\mathrm{Catfib}^{W}_{\deleted{\unsure{W}}}$, \deleted{\unsure{Relativement}} \inserted{ce qui}
(on le verra) \inserted{\unsure{revient}} \deleted{\unsure{naturellement dans l'optique}} en fait au même
% archive sheet number: 307
